\documentclass[10pt]{article}

\usepackage[a4paper,margin=1in]{geometry}
\usepackage[T1]{fontenc}
\usepackage[utf8]{inputenc}
\usepackage{lmodern}
\usepackage{microtype}
\usepackage{booktabs}
\usepackage{graphicx}
\usepackage{amsmath}
\usepackage{hyperref}
\usepackage{enumitem}

\title{Evaluating EDAM-Terms-Reco: A Ground-Truth-Based Assessment of EDAM Term Recommendation}
\author{First Author \and Second Author}
\date{}

\begin{document}
\maketitle

\begin{abstract}
Annotating bioinformatics resources with EDAM terms improves discoverability and interoperability, but manual curation is costly.
We present an evaluation of EDAM-Terms-Reco, a prototype that recommends EDAM ontology terms from textual descriptions.
Using a curated ground truth, we assess recommendation quality with ranking and multi-label metrics.
Results indicate that the approach is promising for top-ranked suggestions, while recall remains limited for rare or highly specific terms.
We discuss error patterns and practical improvements for future iterations.
\end{abstract}

\section{Introduction}
Semantic annotation is essential for FAIR bioinformatics metadata, yet assigning EDAM terms is often time-consuming and inconsistent.
Automated recommendation systems can assist curators by proposing relevant terms from free text.
EDAM-Terms-Reco is an early prototype designed for this purpose.
In this paper, we provide a first ground-truth-based evaluation and answer the following research question:

\begin{quote}
\textit{How accurately can EDAM-Terms-Reco recommend relevant EDAM terms, and how does quality vary across ranking positions and term categories?}
\end{quote}

\section{Methods}
\subsection{Task Definition}
Given a textual description \(x\), the system outputs a ranked list of EDAM terms \((t_1, t_2, \dots)\).
The ground truth is a set of reference terms \(G_x\), making this a multi-label ranking problem.

\subsection{System Under Evaluation}
We evaluate EDAM-Terms-Reco strategies implemented in our prototype (e.g., V1 and variants), which rely on ontology term resources and vector retrieval.

\subsection{Ground Truth}
The benchmark dataset contains bioinformatics resource descriptions with expert-validated EDAM annotations.
We use strict matching at EDAM term ID level.
(Optionally, a relaxed hierarchical analysis can be added in future work.)

\section{Experimental Setup}
\subsection{Data Split and Protocol}
We use a fixed test set (or cross-validation, depending on available data) and keep preprocessing consistent across runs.
Hyperparameters are fixed before final evaluation.

\subsection{Metrics}
We report:
\begin{itemize}[leftmargin=*]
    \item Precision@k and Recall@k (\(k \in \{1,3,5\}\))
    \item F1@k
    \item Mean Average Precision (MAP)
    \item nDCG@k
    \item Exact Match Ratio (strict set equality)
    \item Coverage (fraction of instances with at least one relevant recommendation)
\end{itemize}

\section{Results}
Table~\ref{tab:main-results} summarizes the main performance indicators.
Figure~\ref{fig:pipeline} illustrates the evaluation pipeline.

\begin{table}[h]
\centering
\caption{Main evaluation results on the test set.}
\label{tab:main-results}
\begin{tabular}{lcccccc}
\toprule
Model & P@1 & P@3 & R@3 & MAP & nDCG@5 & Exact Match \\
\midrule
Baseline (lexical) & 0.00 & 0.00 & 0.00 & 0.00 & 0.00 & 0.00 \\
EDAM-Terms-Reco V1 & 0.00 & 0.00 & 0.00 & 0.00 & 0.00 & 0.00 \\
\bottomrule
\end{tabular}
\end{table}

\begin{figure}[h]
\centering
\fbox{\parbox{0.9\linewidth}{
\centering Placeholder for evaluation pipeline figure:
\\ Input description $\rightarrow$ Retrieval/Ranking $\rightarrow$ Top-k EDAM terms $\rightarrow$ Metric computation
}}
\caption{Evaluation workflow.}
\label{fig:pipeline}
\end{figure}

\section{Discussion}
Our initial results suggest that the system can provide useful top-ranked suggestions, supporting curator-in-the-loop workflows.
However, performance drops for rare, ambiguous, or highly specific terms.
Common failure modes include lexical ambiguity and confusion between related EDAM branches (e.g., topic vs operation).
These findings motivate score calibration, better handling of hierarchical structure, and domain-adapted reranking.

\section{Threats to Validity}
Internal validity may be affected by limited benchmark size and annotation noise.
External validity is constrained by domain coverage and ontology version.
Construct validity depends on strict ID matching, which may underestimate semantically close predictions.

\section{Conclusion}
We presented a short, ground-truth-based evaluation of EDAM-Terms-Reco.
The prototype shows encouraging precision at top ranks but still requires improvements in recall and robustness.
Future work will extend benchmarking, include hierarchical metrics, and compare against stronger retrieval baselines.

\section*{Reproducibility Statement}
We provide code, ontology version, benchmark split definition, fixed random seeds, and evaluation scripts to support reproducibility.

\bibliographystyle{plain}
\begin{thebibliography}{9}
\bibitem{edam}
Ison, J. et al. EDAM: an ontology of bioinformatics operations, types of data and identifiers, topics and formats.
\bibitem{fair}
Wilkinson, M. D. et al. The FAIR Guiding Principles for scientific data management and stewardship.
\end{thebibliography}

\end{document}

